\section{One-Time Pad (OTP)}

\subsection{Konsep}

One-Time Pad adalah cipher yang menggunakan kunci acak (random) dengan panjang sama dengan plaintext. Setiap bit/karakter plaintext di-XOR atau dijumlahkan modulo dengan bit/karakter kunci yang sesuai. Kunci hanya digunakan sekali (one-time) \cite{katz2020,stallings2017}.

\[
c_i = p_i \oplus k_i \quad \text{(XOR)} \quad \text{atau} \quad c_i = (p_i + k_i) \bmod 26
\]

\begin{table}[htbp]
  \centering
  \small
  \begin{tabular}{@{}p{3.2cm}p{9.2cm}@{}}
    \toprule
    Sifat & Deskripsi \\
    \midrule
    Keamanan teoretis & Jika kunci benar-benar acak dan dipakai sekali, OTP membuktikan \textit{perfect secrecy} (Shannon) \\
    Kelemahan praktis & Kunci harus sepanjang pesan; distribusi kunci sulit \\
    Penggunaan & Protokol rahasia tingkat tinggi (misal diplomatic), vernam cipher \\
    \bottomrule
  \end{tabular}
  \caption{Sifat One-Time Pad.}
\end{table}

OTP memberikan keamanan sempurna dalam arti teori informasi: ciphertext tidak membocorkan informasi tentang plaintext jika kunci benar-benar acak \cite{rosulek2021}. Namun, pembangkitan dan distribusi kunci yang sepanjang pesan membuat OTP tidak praktis untuk kebanyakan aplikasi. OTP diajarkan sebagai pembanding ideal terhadap cipher klasik lain yang dapat dipecahkan.

\subsection{Contoh Ringkas}

Plaintext: HELLO (7,4,11,11,14). Kunci acak: XMCKL (23,12,2,10,11). Enkripsi mod 26:
\begin{align*}
7+23=4 \to E,\quad 4+12=16 \to Q,\quad 11+2=13 \to N,\\
11+10=21 \to V,\quad 14+11=25 \to Z
\end{align*}
Ciphertext: EQNVZ. Dekripsi: kurangi kunci mod 26.

\begin{catatan}
Menggunakan kunci berulang atau kunci yang dapat diprediksi menghancurkan keamanan OTP. OTP aman \textbf{hanya} jika kunci acak dan dipakai sekali.
\end{catatan}
